\section*{\centering\heiti\zihao{4}附\quad 录}
\addcontentsline{toc}{section}{附录}

{\songti\zihao{-4}
附录补充项目目录、核心命令和提交材料检查表，便于评审复现。

\subsection*{附录A 项目目录说明}

\renewcommand\arraystretch{1.0}
\begin{table}[H]
  \centering
  \caption{项目目录与功能}
  \renewcommand\arraystretch{1.3}
  \setlength{\tabcolsep}{5pt}
  \begin{tabularx}{0.95\textwidth}{>{\raggedright\arraybackslash}p{4.6cm}>{\raggedright\arraybackslash}X}
    \toprule
    路径 & 功能 \\
    \midrule
    \texttt{fraudsim/features.py} & 交易、用户、设备、IP、商户画像合并与训练特征选择。 \\
    \texttt{fraudsim/window\_features.py} & 离线和在线共享的滑动窗口特征计算。 \\
    \texttt{fraudsim/training/train.py} & 多模型训练、阈值校准、指标输出和排行榜更新。 \\
    \texttt{fraudsim/streaming/flink\_job.py} & Flink实时评分作业，输出风险、告警和迟到事件。 \\
    \texttt{fraudsim/api/app.py} & FastAPI模型服务、Dashboard接口、模型热切换、反馈和图挖掘接口。 \\
    \texttt{fraudsim/graph\_mining.py} & 疑似欺诈团伙挖掘、实体风险特征和证据链生成。 \\
    \texttt{fraudsim/dashboard/} & HTML、CSS、JavaScript实现的可视化看板。 \\
    \texttt{scripts/scripts/etl/} & 数据规范化、统一数据集构建、校验和欺诈模式注入。 \\
    \texttt{scripts/benchmark\_api\_latency.py} & 单笔和批量模型API延迟压测。 \\
    \texttt{tests/} & API、特征、窗口、流式微批和图解释单元测试。 \\
    \bottomrule
  \end{tabularx}
\end{table}

\subsection*{附录B 一键演示命令}

\begin{lstlisting}[language=bash]
# 训练模型
python -m fraudsim.training.train --dataset fp_fraudsim_injected --model lightgbm

# 启动基础服务和Flink风险作业
docker compose -p fraudsim up -d --build kafka redis model-api flink-jobmanager flink-taskmanager flink-risk-job

# 创建主题、加载画像、回放交易流
python -m fraudsim.streaming.topics --bootstrap-servers localhost:9094
python -m fraudsim.streaming.load_profiles --dataset fp_fraudsim_injected --redis-url redis://localhost:6379/0
python -m fraudsim.streaming.producer --dataset fp_fraudsim_injected --bootstrap-servers localhost:9094 --rate 100 --limit 1000

# 打开看板
http://localhost:8000/dashboard
\end{lstlisting}

\subsection*{附录C 评分材料检查表}

\begin{table}[H]
  \centering
  \caption{赛题一提交材料检查表}
  \renewcommand\arraystretch{1.3}
  \setlength{\tabcolsep}{5pt}
  \begin{tabularx}{0.95\textwidth}{>{\raggedright\arraybackslash}p{2.5cm}>{\raggedright\arraybackslash}X}
    \toprule
    检查项 & 说明 \\
    \midrule
    技术文档 & 本报告已覆盖任务理解、数据、特征、架构、模型、实验、安全和复现。 \\
    代码工程 & 根目录包含README、Dockerfile、docker-compose、训练、流式、API、Dashboard和测试代码。 \\
    数据与模型 & 大体积目录默认不入库，最终提交前应按赛事要求附带可复现数据包或下载/生成说明。 \\
    指标报告 & 由\texttt{metrics.json}、\texttt{leaderboard.json}和API延迟报告提供。 \\
    演示材料 & Dashboard可展示实时风险、告警、理由码、图团伙和人工反馈闭环。 \\
    安全说明 & 正文第五节给出脱敏、加密、权限、审计、模型安全和容器隔离设计。 \\
    \bottomrule
  \end{tabularx}
\end{table}
}
